% 8.模型改进推广.tex —— 八、模型的改进与推广
\section{模型的改进与推广}

\subsection{模型的改进}
首先补充蒸发速率、潜热及干基质量与体积密度的对应关系，再建立包含水分迁移焓与蒸发耗热的能量平衡。蒸发耗热应以非负蒸发质量率构造热汇，不能未经质量基准和符号核查就直接写成含 $\partial C/\partial t$ 的源项。问题二、三已经通过物性双向耦合，加入潜热的作用是完善能量机制。

其次，问题四把收缩半径 $R(t)$ 当作给定边界，未建立“水分流失量 $\to$ 收缩量”的本构关系。改进方向是补出收缩本构（如体积收缩与含水率变化之间的经验关系），使半径与水分浓度\textbf{双向耦合}，从而在缺少附件 2 数据的新工况下也能外推烘干时长；同时把轴向收缩与端部收缩纳入，建立有限长圆柱的轴对称二维对照模型，检验端部效应和长期径向近似的适用范围；结合实测数据校准物性、传质系数与平衡含水率。边界类型应由界面传递机制和数据确定，不能仅为改善拟合而任意把 Robin 条件替换为固定表面浓度条件。

数值上，可在现有表面加密网格研究的基础上，按终止时间和浓度剖面分别设置精度目标。对达到迭代或子步上限的步，进一步考虑拒绝重算、阻尼或残差控制，并验证这些处理的成本与收益。若将 QS 用于提前预测，应仅使用预测起点之前的数据估计漂移，并用之后未参与拟合的轨道检验其误差。

\subsection{模型的推广}
径向守恒离散与隐式三对角求解可用于其他轴对称传热传质问题，但推广时需重新确认材料关系、几何权重和界面边界。球对称模型需相应调整径向通量和控制体积权重；有限长二维模型的求解结构也随之变化，不能直接照搬一维 Thomas 求解流程。对不同材料与运行条件，均应重新进行参数校准、网格与时间步检查及实验验证。
